% Author-supplied 24 September 2026 manuscript; provenance and local review are recorded separately.
\section{A codimension-two envelope obstruction}
\label{sec:supplied-p031}

The following argument is adapted from the supplied manuscript
\cite{SuppliedP03120260924}.


Let $K$ be a complete nontrivially valued non-Archimedean field. For a finite upper semicontinuous function $H$ on a good strict $K$-analytic space $X$, consider the upper semicontinuous regularization of the supremum of all continuous locally Fubini--Study functions bounded above by $H$. The classical envelope question asks for geometric hypotheses under which this regularization is classically plurisubharmonic, including its analytic-curve tests after extension of the ground field. One proposed weakening of normality or geometric unibranchness everywhere is to impose either condition only in codimension one. The following construction rules out that weakening, even for a geometrically integral surface with a seminormal singularity.

Write $B=K\langle u,v\rangle$, and set
\[
 A=\{f\in B:f(0,0)=f(1,0)\},\qquad X=\mathcal M(A).
\]
The normalization of $X$ is the closed bidisc $\mathcal M(B)$; its map $\nu$ identifies the two rational points $p=(0,0)$ and $q=(1,0)$ to one point $m\in X$.

\begin{theorem}\label{thm:supplied-p031-codimension-two-envelope}
For every complete nontrivially valued non-Archimedean field $K$, the space $X$ is a good strict geometrically integral and geometrically seminormal analytic surface, smooth away from the codimension-two point $m$. There is a finite upper semicontinuous obstacle $H$ on $X$ such that the upper semicontinuous regularization of its full continuous locally Fubini--Study minorant envelope equals $H$, but $H$ is not classically plurisubharmonic. The same envelope identity and the same failure of an analytic-curve test hold after every complete extension $L/K$.
\end{theorem}

The result concerns the sufficiency of codimension-one hypotheses. It leaves open the envelope question under normality or geometric unibranchness at every point. The obstruction is the compatibility of the two branches over $m$: continuity forces every minorant to take the smaller branch value at $m$, whereas upper regularization restores the larger value and creates a positive spike on the other branch.

\subsection{The pinched bidisc}\label{sec:supplied-p031-pinched-bidisc}

Put $s=u(u-1)$, $t=us$, and $w=uv$, and regard $C=K\langle s,v\rangle$ as a subalgebra of $B$. The following explicit decomposition controls the topology as well as the algebra.

\begin{lemma}\label{lem:supplied-p031-banach-decomposition}
The inclusion $C\subset B$ is isometric, and $B=C\oplus uC$ with the maximum norm. Moreover,
\[
 A=C\oplus u(s,v)C=K\langle s,v,t,w\rangle,
\]
where the last equality is an equality of strict affinoid algebras with the induced norm.
\end{lemma}

\begin{proof}
The relation $u^2=u+s$ reduces every polynomial in $u,v$ to a unique expression $F(s,v)+uG(s,v)$. The reduction does not increase the Gauss norm, because the recursion has integral coefficients. Conversely, let $F$ and $G$ be Tate series and suppose their maximum Gauss norm is nonzero. Scale by a coefficient of that norm and reduce the proposed image modulo the maximal ideal of the valuation ring. The homomorphism
\[
 \widetilde K[S,V]\oplus U\widetilde K[S,V]\longrightarrow\widetilde K[U,V],\qquad S\longmapsto U^2-U,
\]
is injective: the leading $U$-degrees of nonzero terms from its two summands have opposite parity. The image therefore has the original maximum norm. Applying this argument to the polynomial truncations and using completeness gives the asserted isometric decomposition for Tate series, including the isometry of $C$.

The equality $f(p)=f(q)$ for $f=F(s,v)+uG(s,v)$ is equivalent to $G(0,0)=0$. Every such $G$ has a decomposition $G=sG_1+vG_2$ with $\max(\|G_1\|,\|G_2\|)\leq\|G\|$, obtained by splitting off its monomials divisible by $s$ and the remaining monomials divisible by $v$. Thus $f=F+tG_1+wG_2$. This expression supplies a bounded quotient map from $K\langle S,V,T,W\rangle$ to $A$ whose quotient norm is the induced norm.
\end{proof}

\begin{proposition}\label{prop:supplied-p031-normalization-base-change}
The inclusion $A\subset B$ is a finite birational normalization. The associated analytic map $\nu:\mathcal M(B)\to X$ is an isomorphism away from $m=\{s=v=t=w=0\}$ and has precisely the two points $p,q$ above $m$. The point $m$ is $K$-rational of codimension two. For every complete extension $L/K$,
\[
 A\widehat\otimes_K L=\{f\in L\langle u,v\rangle:f(0,0)=f(1,0)\},
\]
and all these assertions hold over $L$. In particular, $X$ is geometrically integral and geometrically seminormal, and is smooth outside $m$.
\end{proposition}

\begin{proof}
By \cref{lem:supplied-p031-banach-decomposition}, $B=A[u]$ and $u^2-u=s$. Also $u=t/s$ in the fraction field of $A$, so $A\subset B$ is finite and birational. The Tate algebra $B$ is factorial and hence normal \cite[Theorem~5.4]{JonssonHorawa20}; consequently it is the integral closure of $A$.

On the locus $s\neq0$, the identity $u=t/s$ defines an analytic inverse to $\nu$. On $v\neq0$ the inverse is $u=w/v$, and the two expressions agree on their overlap. The relations
\[
 t^2-st-s^3=0,\qquad w^2-vw-sv^2=0,\qquad tv=sw
\]
show that a point with $s=v=0$ also has $t=w=0$. Thus the complement of these two inverse charts is the single rational point $m$. Since $B/(s,v)=K[u]/(u(u-1))\simeq K\times K$, the fiber of $m$ consists of $p$ and $q$. The affinoid spectra are compact Hausdorff, so the continuous surjection $\nu$ is a quotient map.

The finite extension $C\subset A$ and the dimension of $C$ give $\dim A=2$. The ideal $(v,w)=vB$ in $A$ is prime: its quotient is $\{f\in K\langle u\rangle:f(0)=f(1)\}$, a domain. It lies strictly between $(0)$ and the maximal ideal $m$, so $m$ has height two. The isomorphism away from $m$ identifies that locus with an open subset of the smooth bidisc.

The decomposition $A=C\oplus u(s,v)C$ is a split Banach-space decomposition. Completed extension of scalars therefore gives the same equalizer description over $L$, and the preceding normalization and fiber arguments apply verbatim to $L$. In particular $A\widehat\otimes_K L$ is a domain. For seminormality, let $b$ in its fraction field satisfy $b^2,b^3\in A\widehat\otimes_K L$. It is integral and therefore belongs to the normalizing algebra $L\langle u,v\rangle$. If $\alpha=b(p)$ and $\beta=b(q)$, then $\alpha^2=\beta^2$ and $\alpha^3=\beta^3$, forcing $\alpha=\beta$ in every characteristic. Hence $b$ lies in the equalizer. This is the square-and-cube criterion \cite[Definition~29.48.1, \href{https://stacks.math.columbia.edu/tag/0EUL}{Tag~0EUL}]{StacksProject}, and it proves seminormality after every complete extension. The two distinct points $p,q$ above $m$ also show that this singularity is not geometrically unibranch.
\end{proof}

\subsection{The full minorant envelope}\label{sec:supplied-p031-envelope}

Choose $\pi\in K$ with $0<|\pi|<1$ and put $c=-\log|\pi|>0$. On $\mathcal M(B)$ consider the continuous function
\[
 h=\max\{0,c+\log|u-1|\},
\]
with the value $0$ at $u=1$. It takes values in $[0,c]$, equals $c$ where $|u|<1$, and equals $0$ where $|u-1|\leq|\pi|$. In particular $h(p)=c$ and $h(q)=0$. Define $H(m)=c$ and $H(x)=h(\nu^{-1}(x))$ for $x\neq m$. The function $H$ is upper semicontinuous: it is continuous off $m$ and its value at $m$ is its global upper bound.

Let $\mathcal F_H$ be the full family of continuous locally Fubini--Study functions $\varphi$ on $X$ with $\varphi\leq H$, and let $g=\sup_{\varphi\in\mathcal F_H}\varphi$. The constant zero function belongs to $\mathcal F_H$.

\begin{proposition}\label{prop:supplied-p031-envelope-identity}
For this obstacle, $g(m)=0$, $g=H$ on $X\setminus\{m\}$, and $g^*=H$ on $X$. The same identities hold for the full family of continuous semipositive real piecewise-linear minorants whenever that class contains global Fubini--Study functions.
\end{proposition}

\begin{proof}
For $n\geq1$ define $f_n=(u-1)^n s$ and $k_n=(u-1)^n v$. Both lie in $A$ because they vanish at $p$ and $q$. The function
\begin{equation}\label{eq:supplied-p031-fs-minorants}
 \psi_n=\frac1n\log\max\bigl\{1,|\pi|^{-n}|f_n|,|\pi|^{-n}|k_n|\bigr\}
\end{equation}
is globally Fubini--Study: its three sections of the trivial bundle are $1,\pi^{-n}f_n,\pi^{-n}k_n$, and the first section removes any base locus. Set $R=\max\{|s|,|v|\}\leq1$ on $\mathcal M(B)$. Then
\[
 \psi_n\circ\nu=\max\left\{0,c+\log|u-1|+\frac1n\log R\right\},
\]
interpreted as zero at $R=0$ or $u=1$. Thus $0\leq\psi_n\leq H$, including at $m$ where $\psi_n(m)=0$. Away from $p,q$ one has $R>0$, so $\psi_n(x)\to H(x)$ for each $x\neq m$. It follows that $g=H$ away from $m$.

Every continuous minorant $\varphi\leq H$ satisfies $\varphi(m)\leq0$: the points $\nu(1,\pi^j)$ tend to $m$, avoid it, and have obstacle value zero. Since the zero function is a minorant, $g(m)=0$. Conversely, $\nu(0,\pi^j)\to m$ and $g(\nu(0,\pi^j))=H(\nu(0,\pi^j))=c$. Therefore $g^*(m)\geq c$. The inequality $g\leq H$ and upper semicontinuity of $H$ give $g^*\leq H$, proving $g^*=H$. This argument only uses continuity and the presence of the functions $0,\psi_n$ in the chosen full minorant class, so it applies to the stated semipositive real piecewise-linear class as well.
\end{proof}

\begin{proof}[of \Cref{thm:supplied-p031-codimension-two-envelope}]
For the geometry apply \cref{prop:supplied-p031-normalization-base-change}; for the envelope apply \cref{prop:supplied-p031-envelope-identity}. To test positivity, let $\gamma:\mathcal M(K\langle z\rangle)\to X$ be the analytic map induced by $u=1$, $v=z$ on the normalization. Its pullback is
\[
 (g^*\circ\gamma)(z)=
 \begin{cases}
 c,&z=0,\\
 0,&z\neq0.
 \end{cases}
\]
On the closed subdisc $|z|\leq|\pi|$, the origin is a relative-interior point, and the function is zero at its Gauss boundary point. Comparison with the harmonic function zero would bound its value at the origin by zero, contrary to $c>0$. Thus $g^*$ fails the analytic-curve test for classical plurisubharmonicity at an interior point. The description of $A$ in \cref{prop:supplied-p031-normalization-base-change} shows that the same formulas define the scalar extensions of $H$, $\psi_n$, and $\gamma$. The two sequences and the disc comparison then prove the claims over every complete extension $L/K$.
\end{proof}
